\documentclass[11pt]{amsart}

\usepackage[T1]{fontenc}
\usepackage{lmodern}
\usepackage{microtype}
\usepackage{mathtools,amssymb,amsthm}
\usepackage{array}
\usepackage[hidelinks]{hyperref}

\hypersetup{
  pdftitle={The hexagonal prism under a closed-neighbourhood one-round rule:
    fractional eternal domination number 46/13}
}

\newtheorem{theorem}{Theorem}
\newtheorem{proposition}[theorem]{Proposition}
\newtheorem{lemma}[theorem]{Lemma}
\newtheorem{corollary}[theorem]{Corollary}
\theoremstyle{remark}
\newtheorem*{remark}{Remark}
\newtheorem*{problem}{Problem 7.5 of \cite{FED}}

\newcommand{\fed}{\gamma_f^{\infty}}
\newcommand{\Z}{\mathbb{Z}}
\newcommand{\Cay}{\operatorname{Cay}}
\newcommand{\half}{\tfrac12}

\title[The hexagonal prism under a closed-neighbourhood rule]
{The Hexagonal Prism under a Closed-Neighbourhood One-Round Rule:
 Fractional Eternal Domination Number $46/13$}

\date{}

\subjclass[2020]{05C69, 05C57, 05C25, 90C05}
\keywords{fractional eternal domination, cubic abelian Cayley graph, prism,
linear programming duality, exact value}

\begin{document}

\begin{abstract}
For the hexagonal prism $C_6\mathbin{\square}K_2$, and for the
\emph{closed-neighbourhood} reading of the defender's one-round rule set out in
Section~\ref{sec:intro},
\[
  \fed(C_6\mathbin{\square}K_2)=\frac{46}{13}=3.\overline{538461}\,,
\]
and the infimum is attained. That reading is ours and not a quotation: the
source states the rule in a form we read as leaving unmoved weight in place, and
both certificates below encode that reading, so if the intended rule is stricter
they concern a different game and the value above is not the value asked for.
This graph is the cell $k=1$ of the first family of Problem 7.5 of Devvrit,
Krim-Yee, Kumar, MacGillivray, Seamone, Virgile and Xu, who asked for the exact
value of the fractional eternal domination number $\fed(G)$ when $G$ is a cubic
abelian Cayley graph of the two exceptional types $C_{4k+2}\mathbin{\square}K_2$
($k\ne2$) and $\Cay(\Z_{8k},\{\pm1,4k\})$, and who prove for this cell only
$\tfrac72<\fed(C_6\mathbin{\square}K_2)\le 4$. The upper bound here is an
explicit $12$-member family of rational weight functions, printed in full; the
lower bound is a $9$-node attack tree together with a dual certificate in exact
rationals, also printed in full, so both halves can be checked against the
printed objects alone --- the $493$ dual column inequalities being tedious
rather than difficult. A closing remark records that, for the same reading of
the rule, the last step of the published proof of $\tfrac72<\fed$ is not
justified; the published inequality is nevertheless true, since
$\tfrac72<\tfrac{46}{13}$. No other cell of the
authors' question is settled here, and we do not settle which rule they
intended.
\end{abstract}

\maketitle

\section{Introduction}\label{sec:intro}

We use the definitions of Devvrit et al.\ \cite{FED}. A
\emph{fractional dominating function} of a graph $G$ is a map
$w\colon V(G)\to\mathbb{R}_{\ge0}$ with $\sum_{x\in N[v]}w(x)\ge1$ for every
$v\in V(G)$, and $G$ is \emph{$S$-fractionally dominated} if it admits a
fractional dominating function of total weight at most $S$. In the eternal
game, weights are assigned so that $S$-fractional domination is maintained
under vertex attacks: writing $w_i$ for the weight function at time step $i$
and $m_{xy,i}$ for the weight moved from $x$ to $y\in N(x)$ in round $i$, the
defender may move simultaneously from as many vertices as desired subject to
$\sum_{y\in N(x)}m_{xy,i}\le w_i(x)$, and the resulting $w_{i+1}$ must again
$S$-fractionally dominate $G$ and satisfy $w_{i+1}(v_i)\ge1$, where $v_i$ is
the vertex attacked in round $i$. The \emph{fractional eternal domination
number} $\fed(G)$ is the infimum of those $S$ for which $G$ can be eternally
$S$-fractionally dominated.

Two features of this definition are used below, and the first of them is our
reading of the source rather than a quotation from it. The budget constraint we
have quoted is stated in \cite{FED} as
$\sum_{y\in N(x)}m_{xy,i}\le w_i(x)$, an open-neighbourhood sum bounded by the
weight at $x$; we read it as leaving unmoved weight where it is, so that a
legal round transports weight along \emph{closed} neighbourhoods with budget
$w_i(x)$ at $x$. Everything below --- the family of
Section~\ref{sec:upper} and the linear program of Section~\ref{sec:lower}
alike --- encodes that reading, and if the intended rule is stricter (for
instance if weight must leave $x$ along open neighbourhoods) then both encode a
different game and the value below is not the value asked for. We have not
found in \cite{FED} an update equation for $w_{i+1}$ that decides between the
two readings, and we do not decide it here; that is why
Theorem~\ref{thm:main} below is stated for the closed-neighbourhood rule rather
than as a determination of the cell the authors asked about. The accompanying
program encodes the same reading, so no check in it can discriminate between
the two. Second, in the
fractional model of \cite{FED} the attacker is unrestricted, in particular it
may attack a vertex that already carries weight.

Devvrit et al.\ prove that a cubic abelian Cayley graph $G$ satisfies
$\fed(G)<\gamma(G)$ if and only if $G\cong C_{4k+2}\mathbin{\square}K_2$ or
$G\cong\Cay(\Z_{8k},\{\pm1,4k\})$ for some positive integer $k$, and compute
$\fed$ exactly for all other cubic abelian Cayley graphs. For the two
exceptional families they ask:

\begin{problem}
Let $k$ be a positive integer. Determine the exact value of $\fed(G)$ when $G$
is isomorphic to
\begin{itemize}
\item $C_{4k+2}\mathbin{\square}K_2$ for $k\ne2$;
\item $\Cay(\Z_{8k},\{\pm1,4k\})$.
\end{itemize}
\end{problem}

\noindent
Ordering the cells of this question by $|V|$, the value is already a theorem of
\cite{FED} at $|V|=8$, namely $\fed(\Cay(\Z_8,\{\pm1,4\}))=\tfrac83$, and at
$|V|=20$, namely $\fed(C_{10}\mathbin{\square}K_2)=\tfrac{28}{5}$ (the cell
$k=2$ excluded by the first bullet). The smallest cell it leaves open is
therefore $k=1$ of the first bullet, the hexagonal prism
$C_6\mathbin{\square}K_2$ on $12$ vertices, where \cite{FED} proves
\[
  \tfrac72<\fed(C_6\mathbin{\square}K_2)\le4 .
\]
The upper bound there is derived from the general estimate
$\fed(G)\le(n+\kappa)/(\kappa+1)$ of \cite{FED} for a graph of connectivity
$\kappa$, which at $n=12$, $\kappa=3$ gives $\tfrac{15}{4}$; so \cite{FED} in
fact also yields $\fed(C_6\mathbin{\square}K_2)\le\tfrac{15}{4}=3.75$ here.

\begin{theorem}\label{thm:main}
For the closed-neighbourhood one-round rule read off above,
$\displaystyle \fed(C_6\mathbin{\square}K_2)=\frac{46}{13}$, and the infimum is
attained: an eternal strategy of total weight exactly $\tfrac{46}{13}$ exists.
\end{theorem}

The value is consistent with the published statements about this cell:
$\tfrac72<\tfrac{46}{13}<\tfrac{15}{4}<4$, with
$\tfrac{46}{13}-\tfrac72=\tfrac1{26}$; and
$\tfrac{46}{13}-\gamma_f(C_6\mathbin{\square}K_2)=\tfrac{46}{13}-3=\tfrac7{13}<1$,
as the corollary $\fed-\gamma_f<1$ of \cite{FED} for cubic abelian Cayley
graphs demands. (Here $\gamma_f(C_6\mathbin{\square}K_2)=3$: the uniform weight
$\tfrac14$ is fractionally dominating with total $3$, and summing the twelve
closed neighbourhood inequalities gives $4S\ge12$.)
Section~\ref{sec:upper} proves the upper bound and Section~\ref{sec:lower} the
lower bound, closing with a remark on the published proof of
$\tfrac72<\fed$; Section~\ref{sec:open} states what is \emph{not} settled.

\medskip
\noindent\textbf{Notation.} Throughout, $G=C_6\mathbin{\square}K_2$ with vertex
set $V=\Z_6\times\{0,1\}$, edges $(i,j)\sim(i\pm1,j)$ and $(i,0)\sim(i,1)$.
We also write $v=i+6j$, so that $V=\{0,\dots,11\}$ with
\begin{center}\footnotesize
\begin{tabular}{llll}
$N[0]=\{0,1,5,6\}$   & $N[1]=\{0,1,2,7\}$   & $N[2]=\{1,2,3,8\}$    & $N[3]=\{2,3,4,9\}$\\
$N[4]=\{3,4,5,10\}$  & $N[5]=\{0,4,5,11\}$  & $N[6]=\{0,6,7,11\}$   & $N[7]=\{1,6,7,8\}$\\
$N[8]=\{2,7,8,9\}$   & $N[9]=\{3,8,9,10\}$  & $N[10]=\{4,9,10,11\}$ & $N[11]=\{5,6,10,11\}$.
\end{tabular}
\end{center}
Thus $G$ is $3$-regular with $18$ edges, $3$-connected, bipartite and of girth
$4$.

\section{The upper bound}\label{sec:upper}

\begin{lemma}\label{lem:frame}
Let $\mathcal F=\{f_v : v\in V\}$ be a family of fractional dominating
functions of $G$, indexed by $V$, such that
\begin{enumerate}
\item[(i)] each $f_v$ has total weight $S$;
\item[(ii)] $f_v(v)=1$ for every $v$;
\item[(iii)] for every ordered pair $(u,v)\in V\times V$ there is a
  transportation plan $x\colon\{(a,b):b\in N[a]\}\to\mathbb{R}_{\ge0}$ with
  $\sum_{b\in N[a]}x(a,b)=f_u(a)$ for all $a$ and
  $\sum_{a:\,b\in N[a]}x(a,b)=f_v(b)$ for all $b$.
\end{enumerate}
Then $\fed(G)\le S$, and the infimum is attained.
\end{lemma}

\begin{proof}
Play from $f_{v_0}$ for any $v_0$, and respond to an attack at $v$, from the
current configuration $f_u$, by moving to $f_v$; (iii) says the move is legal
for the one-round rule, since weight sent from $a$ to $a$ itself is weight left
in place, (i) says the total is unchanged, $f_v$ is fractionally dominating,
and (ii) says $f_v(v)=1\ge1$. The strategy therefore survives every infinite
sequence of attacks at total weight exactly $S$.
\end{proof}

For $(m,t)\in V$ put $f_{(m,t)}(i,j)=h_{j\oplus t}\bigl((i-m)\bmod6\bigr)$,
where $\oplus$ is addition mod $2$ and
\[
\begin{aligned}
  h_0&:\ 0\mapsto1,\quad \pm1\mapsto\tfrac1{13},\quad
        \pm2\mapsto\tfrac1{13},\quad 3\mapsto\tfrac5{13};\\
  h_1&:\ 0\mapsto\tfrac6{13},\quad \pm1\mapsto0,\quad
        \pm2\mapsto\tfrac6{13},\quad 3\mapsto\tfrac6{13}.
\end{aligned}
\]
Both profiles are even, so each $f_{(m,t)}$ is invariant under the stabiliser
$\{\mathrm{id},\,i\mapsto-i\}$ of $(0,0)$ and the twelve translates are well
defined. Written out, the seed $f=f_{(0,0)}$ is
\[
  f=\Bigl(1,\tfrac1{13},\tfrac1{13},\tfrac5{13},\tfrac1{13},\tfrac1{13}
   \,\Big|\,\tfrac6{13},0,\tfrac6{13},\tfrac6{13},\tfrac6{13},0\Bigr)
\]
in the order $v=0,\dots,11$, of total weight
$1+\tfrac9{13}+\tfrac{24}{13}=\tfrac{46}{13}$. Its twelve closed neighbourhood
sums are
\[
  \tfrac{21}{13},\ \tfrac{15}{13},\ 1,\ 1,\ 1,\ \tfrac{15}{13}
  \ \Big|\
  \tfrac{19}{13},\ 1,\ 1,\ \tfrac{23}{13},\ 1,\ 1,
\]
all at least $1$; translation carries this to every
member, so (i) and (ii) of Lemma~\ref{lem:frame} hold with
$S=\tfrac{46}{13}$.

For (iii), here is one plan in full, for $f_{(0,0)}\to f_{(3,0)}$; the notation
$a\!\to\!b=q$ means that $a$ ships $q$ to $b$, and $b\in N[a]$ in every case:
\begin{center}\ttfamily\small
0->0=5/13\quad 0->1=1/13\quad 0->5=1/13\quad 0->6=6/13\quad 1->2=1/13\\
2->3=1/13\quad 3->3=5/13\quad 4->3=1/13\quad 5->4=1/13\quad 6->7=6/13\\
8->9=6/13\quad 9->3=6/13\quad 10->11=6/13
\end{center}
Every vertex ships exactly its own weight and every vertex receives exactly the
target weight,
\[
\begin{aligned}
 \text{ships}&:\ \bigl(1,\tfrac1{13},\tfrac1{13},\tfrac5{13},\tfrac1{13},
   \tfrac1{13}\mid\tfrac6{13},0,\tfrac6{13},\tfrac6{13},\tfrac6{13},0\bigr),\\
 \text{receives}&:\ \bigl(\tfrac5{13},\tfrac1{13},\tfrac1{13},1,\tfrac1{13},
   \tfrac1{13}\mid\tfrac6{13},\tfrac6{13},0,\tfrac6{13},0,\tfrac6{13}\bigr),
\end{aligned}
\]
the second of which is $f_{(3,0)}$. The total shipped
is $\tfrac{46}{13}$, and the attacked vertex $3$ receives
$\tfrac1{13}+\tfrac5{13}+\tfrac1{13}+\tfrac6{13}=1$. All $144$ ordered pairs
of the family are feasible for (iii); each is a transportation problem with
integer data after scaling by $13$, and the accompanying program decides all
$144$ by exact integer maximum flow.

By Lemma~\ref{lem:frame}, $\fed(C_6\mathbin{\square}K_2)\le\tfrac{46}{13}$.

\section{The lower bound}\label{sec:lower}

\subsection{A finite attack tree is a relaxation}
Fix a finite rooted tree $T$ with nodes $0,1,\dots,N$, root $0$, parent map
$p$, and an attacked vertex $a_k\in V$ for each $k\ge1$. Consider

\begin{equation}\label{eq:lp}
\begin{aligned}
\text{minimise } & S\\
\text{subject to } &\textstyle\sum_{v\in V}w_k(v)=S && (0\le k\le N)\\
 &\textstyle\sum_{u\in N[v]}w_k(u)\ge1 && (0\le k\le N,\ v\in V)\\
 & w_k(a_k)\ge1 && (1\le k\le N)\\
 &\textstyle\sum_{y\in N[u]}x_k(u,y)=w_{p(k)}(u) && (1\le k\le N,\ u\in V)\\
 &\textstyle\sum_{u:\,y\in N[u]}x_k(u,y)=w_k(y) && (1\le k\le N,\ y\in V)\\
 & S\ge0,\ w_k(v)\ge0,\ x_k(u,y)\ge0. &&
\end{aligned}
\end{equation}

\begin{lemma}\label{lem:relax}
If $L$ is the optimum of \eqref{eq:lp} then $L\le\fed(G)$.
\end{lemma}

\begin{proof}
An eternal defender surviving at total weight $S$ survives, in particular, the
attacks of $T$: taking $w_0$ to be its initial configuration, $w_k$ the
configuration reached along the path from the root to $k$, and $x_k$ the moves
it makes in that round, every constraint of \eqref{eq:lp} is one of the game's
own requirements. So $S$ is feasible for \eqref{eq:lp} and $L\le S$; take the
infimum over $S$.
\end{proof}

Three remarks make Lemma~\ref{lem:relax} robust rather than delicate. The root
configuration $w_0$ is \emph{free}, so no uniqueness or symmetry lemma about
optimal initial configurations is needed. Any subtree of a legal attacker is a
legal attacker, so pruning $T$ can only decrease $L$ --- it can never
invalidate the bound. And an attacker that plays one fixed first attack and
only then branches is weaker than one that branches immediately, so using such
a tree is safe.

\subsection{The tree}
Take $N=8$ with
\[
\begin{array}{l@{\qquad}l}
 p(1)=0, & a_1=(0,0)\\
 p(2)=1,\ p(4)=1,\ p(7)=1, & a_2=(2,1),\ a_4=(3,1),\ a_7=(4,1)\\
 p(3)=2,\ p(5)=p(6)=4,\ p(8)=7, & a_3=a_6=a_8=(5,1),\ a_5=(1,1).
\end{array}
\]
The second attacks $\{(2,1),(3,1),(4,1)\}$ are the three vertices of the other
hexagon at distances $3,4,3$ from the first attack. For this tree
\eqref{eq:lp} has $493$ variables ($S$; the $9\times12$ weights $w_k(v)$; and
$8\times48$ flows $x_k(u,y)$, one per node $k\ge1$ and per ordered pair
$(u,y)$ with $y\in N[u]$), $201$ equalities (nine total-weight rows, then
$8\times(12+12)$ flow rows) and $116$ inequalities ($9\times12$ domination
rows and $8$ attack rows).

\begin{proposition}\label{prop:lower}
The optimum of \eqref{eq:lp} for this tree is exactly $\tfrac{46}{13}$;
consequently $\fed(C_6\mathbin{\square}K_2)\ge\tfrac{46}{13}$.
\end{proposition}

\begin{proof}
Only ``$\ge$'' is needed for the theorem, and it follows from weak duality once
a dual feasible point of objective $\tfrac{46}{13}$ is exhibited. Let
$\alpha_k\ge0$ be the multiplier of the row $w_k(a_k)\ge1$, let
$\beta_{k,v}\ge0$ be the multiplier of the row
$\sum_{u\in N[v]}w_k(u)\ge1$, and let $\lambda_k$, $\mu^{\mathrm{out}}_{k,u}$,
$\mu^{\mathrm{in}}_{k,y}$ be the free multipliers of the total-weight,
out-flow and in-flow equalities respectively. Dual feasibility is the system of
$493$ column inequalities of \eqref{eq:lp}, one per primal variable, and the
dual objective is $\sum_k\alpha_k+\sum_{k,v}\beta_{k,v}$.

Take $\alpha_1=\tfrac4{13}$, $\alpha_2=\tfrac2{13}$, $\alpha_4=\tfrac4{13}$,
$\alpha_5=\tfrac3{13}$, $\alpha_6=\tfrac2{13}$, $\alpha_7=\tfrac2{13}$,
$\alpha_3=\alpha_8=0$, and let $13\,\beta_{k,v}$ be given by
\[
\small
\begin{array}{l@{\ }l}
 k=1: & (1,2)\mapsto1,\ (1,3)\mapsto1,\ (1,4)\mapsto1,\ (1,7)\mapsto2,\\
      & (1,8)\mapsto1,\ (1,10)\mapsto1,\ (1,11)\mapsto2\\
 k=2: & (2,1)\mapsto1,\ (2,3)\mapsto1,\ (2,11)\mapsto1\\
 k=4: & (4,1)\mapsto1,\ (4,2)\mapsto2,\ (4,4)\mapsto2,\ (4,5)\mapsto1,\\
      & (4,6)\mapsto1,\ (4,11)\mapsto1\\
 k=5: & (5,0)\mapsto1,\ (5,2)\mapsto1,\ (5,10)\mapsto1\\
 k=6: & (6,0)\mapsto1,\ (6,4)\mapsto1,\ (6,8)\mapsto1\\
 k=7: & (7,3)\mapsto1,\ (7,5)\mapsto1,\ (7,7)\mapsto1,
\end{array}
\]
all other $\beta_{k,v}$ being $0$. There are $31$ nonzero multipliers on the
$\ge$ rows; the attack part sums to
$(4+2+4+3+2+2)/13=\tfrac{17}{13}$ and the domination part to
$(9+3+8+3+3+3)/13=\tfrac{29}{13}$, so the dual objective is exactly
$\tfrac{17}{13}+\tfrac{29}{13}=\tfrac{46}{13}$.

The free multipliers are $\lambda_3=-\tfrac7{13}$,
$\lambda_4=-\tfrac1{13}$, $\lambda_5=-\tfrac3{13}$,
$\lambda_6=-\tfrac2{13}$, all other $\lambda_k=0$, together with the
following, where each row lists $13\mu^{\mathrm{out}}_{k,u}$ for
$u=0,1,\dots,11$ and then $13\mu^{\mathrm{in}}_{k,y}$ for $y=0,1,\dots,11$;
all multipliers not listed --- in particular every flow multiplier of nodes
$1$ and $8$ --- are $0$:
\begin{center}\ttfamily\small
\begin{tabular}{@{}r@{\ \ }l@{}}
k=2 out: &  6\ \ 5\ \ 5\ \ 5\ \ 6\ \ 6\ \ 6\ \ 5\ \ 5\ \ 5\ \ 6\ \ 6\\
k=2 in:  & -6\ -6\ -5\ -6\ -6\ -6\ -6\ -6\ -5\ -6\ -6\ -6\\
k=3 out: &  7\ \ 7\ \ 7\ \ 7\ \ 7\ \ 7\ \ 7\ \ 7\ \ 7\ \ 7\ \ 7\ \ 7\\
k=3 in:  & -7\ -7\ -7\ -7\ -7\ -7\ -7\ -7\ -7\ -7\ -7\ -7\\
k=4 out: & -1\ -1\ -1\ \ 0\ -1\ -1\ -1\ -1\ \ 0\ \ 0\ \ 0\ -1\\
k=4 in:  &  1\ \ 1\ \ 0\ \ 0\ \ 0\ \ 1\ \ 1\ \ 0\ \ 0\ \ 0\ \ 0\ \ 0\\
k=5 out: &  1\ \ 0\ \ 1\ \ 2\ \ 2\ \ 2\ \ 0\ \ 0\ \ 0\ \ 2\ \ 2\ \ 2\\
k=5 in:  & -2\ -1\ -2\ -2\ -2\ -2\ -2\ \ 0\ -2\ -2\ -2\ -2\\
k=6 out: &  0\ \ 1\ \ 1\ \ 1\ \ 0\ \ 0\ \ 0\ \ 1\ \ 1\ \ 1\ \ 0\ \ 0\\
k=6 in:  & -1\ -1\ -1\ -1\ -1\ \ 0\ -1\ -1\ -1\ -1\ -1\ \ 0\\
k=7 out: & -1\ -1\ -1\ -2\ -2\ -2\ -1\ -1\ -1\ -2\ -2\ -2\\
k=7 in:  &  1\ \ 1\ \ 1\ \ 1\ \ 2\ \ 1\ \ 1\ \ 1\ \ 1\ \ 1\ \ 2\ \ 1
\end{tabular}
\end{center}
This is $123$ nonzero equality multipliers, every one of denominator $13$.
Substituting into the $493$ column inequalities of \eqref{eq:lp} verifies them
all in exact rational arithmetic; the accompanying program does exactly that,
having first rebuilt \eqref{eq:lp} from the description above. Weak duality now
gives $L\ge\tfrac{46}{13}$, and Section~\ref{sec:upper} together with
Lemma~\ref{lem:relax} gives $L\le\fed(G)\le\tfrac{46}{13}$, so
$L=\tfrac{46}{13}$ as well.
\end{proof}

Theorem~\ref{thm:main} follows from Section~\ref{sec:upper} and
Proposition~\ref{prop:lower}.

\begin{remark}
The inequality $\tfrac72<\fed(C_6\mathbin{\square}K_2)$ is true, by
Theorem~\ref{thm:main}, but its published proof ends with a step that is not
justified for the reading of the rule used here; we record this because a reader
reconstructing the lower bound from \cite{FED}
will otherwise reconstruct it. In the labelling of the relevant figure of
\cite{FED} --- inner cycle $a\,b\,c\,d\,e\,f$, outer cycle
$g\,h\,i\,j\,k\,l$, with $a\sim g$, $b\sim h$, \ldots, $f\sim l$ --- the
argument assumes a first attack at $a$ and derives a family of configurations
parameterised by $x\in[0,\half]$, namely $a=1$, $b=f=h=l=0$, $c=e=x$,
$d=j=\half-x$, $g=i=k=\half$. Its last step asserts that, after a response to a
second attack at $j$, ``there must be a weight of $1$ on $j$, $\half$ on $d$
and $0$ on each of the vertices $c,e,h,i,k,l$''. For the closed-neighbourhood
reading used throughout this note that is not forced: start from the $x=0$
member
\[
 a=1,\quad d=g=i=j=k=\tfrac12,\quad b=c=e=f=h=l=0
\]
of total weight $\tfrac72$, and answer an attack at $j$ by the one-round move
\[
\begin{aligned}
 &a\to b=\tfrac12,\quad a\to f=\tfrac12,\quad g\to a=\tfrac14,\\
 &i\to h=\tfrac14,\quad i\to j=\tfrac14,\quad
  k\to l=\tfrac14,\quad k\to j=\tfrac14,
\end{aligned}
\]
with $d$, $j$ and the remaining $\tfrac14$ at $g$ left in place. No vertex
ships more than it holds, the total is conserved, and the resulting
configuration
\[
\begin{aligned}
 &a=\tfrac14,\ b=\tfrac12,\ c=0,\ d=\tfrac12,\ e=0,\ f=\tfrac12,\\
 &g=\tfrac14,\ h=\tfrac14,\ i=0,\ j=1,\ k=0,\ l=\tfrac14
\end{aligned}
\]
puts weight $1$ on the attacked vertex $j$ and has all twelve closed
neighbourhood sums at least $1$. So the attack at $j$ \emph{is} answerable at
total weight $\tfrac72$ from that configuration, and $h=l=\tfrac14\ne0$. The
accompanying program checks this configuration and this response arc by arc.
Whether the step can be justified under a stricter reading of the movement rule
we do not say.
\end{remark}

\section{What is not settled}\label{sec:open}

\begin{enumerate}
\item \textbf{Which game is being played is not settled.} The value
  $\tfrac{46}{13}$ is a value for the closed-neighbourhood one-round rule of
  Section~\ref{sec:intro}, which is our reading of \cite{FED} and not a
  quotation from it. Both certificates above and the accompanying program encode
  that reading, so none of them discriminates between it and a stricter one; for
  a stricter rule nothing here bounds $\fed$, and the cell of the authors'
  question stays open.
\item \textbf{The authors' question is not determined.} At most one cell of it
  is. If the reading above is the intended one, the cells left open after
  Theorem~\ref{thm:main} are
  $\{C_{4k+2}\mathbin{\square}K_2 : k\ge3\}\cup
   \{\Cay(\Z_{8k},\{\pm1,4k\}) : k\ge2\}$: every cell of both families except
  $\tfrac{46}{13}$ and $\tfrac{28}{5}$ at $k=1,2$ of the first and $\tfrac83$ at
  $k=1$ of the second. Nothing here should be read as settling the question as a
  whole.
\item \textbf{No general theorem follows.} The method is cell by cell: an
  explicit finite family of weight functions, and a finite attack tree. It
  gives no construction for general $k$, and it does not address the authors'
  separate questions of whether $\fed$ is always rational, or whether $n$
  fractional dominating functions always suffice --- the latter is what makes a
  bounded search for a family inconclusive rather than a lower bound when it
  finds nothing.
\item \textbf{Only the $9$-node dual is certified here.} Other attack trees were
  solved in the course of the work; none of their optima is quoted or used
  above, and no certificate for any of them is printed. Theorem~\ref{thm:main}
  rests on the $9$-node dual printed above and on nothing else.
\item \textbf{The certificates were found numerically.} A floating-point
  linear programming solver located the witness family and the dual vector.
  That is how they were found and not why they hold: every object above is
  exhibited as exact rationals and every verification --- fractional
  domination, transportation feasibility, the $493$ dual column inequalities,
  the objective --- is carried out in exact integer or rational arithmetic. No
  floating-point value appears anywhere in the chain of reasoning.
\end{enumerate}

\section{Reproduction and verification}

Everything needed is printed above: the family is determined by the profiles
$h_0,h_1$ of Section~\ref{sec:upper}; the lower bound by the attack tree, the
linear program \eqref{eq:lp} with its stated row and variable ordering, and the
$31+123$ multipliers of Proposition~\ref{prop:lower}. Reconfigurability of a
single ordered pair is a transportation problem with integer data after scaling
by $13$, and the plan printed in Section~\ref{sec:upper} can be checked by
hand; the dual objective $\tfrac{17}{13}+\tfrac{29}{13}$ can be re-added by
hand; the $493$ dual column inequalities are the one part that is tedious
rather than difficult.

A program accompanying this note re-derives all of it from the printed objects,
in exact integer and \texttt{Fraction} arithmetic, using only the Python
standard library: it rebuilds $G$ and checks it against the closed
neighbourhoods printed above; rebuilds the family from its profiles and checks
total weight, fractional domination and the centre condition for all $12$
members; checks the printed transportation plan arc by arc; decides all $144$
ordered pairs by exact integer maximum flow; rebuilds \eqref{eq:lp} and
verifies its dimensions $493/201/116$, then verifies all $493$ dual column
inequalities and the dual objective; and re-checks the configuration and the
response of the remark that closes Section~\ref{sec:lower}. What it does not do
is re-run the floating-point linear programs that found the family and the dual
vector; those are re-checked only as exhibited objects. Nor can it bear on the
reading of the movement rule: it encodes the same reading, in the same closed
neighbourhoods.

The output of the recorded run is wider than the note: besides the checks just
listed it records controls, a second family of weight functions of total weight
$\tfrac{46}{13}$, three further transportation plans, an ancillary upper bound
for $\Cay(\Z_{16},\{\pm1,8\})$, and a closing scope note whose phrase ``the
accompanying note'' should not be read as a scope statement about the sections
above. Nothing above depends on any of that extra material, and none of it is
claimed here; what is and is not settled here is stated in
Section~\ref{sec:open}.

\begin{thebibliography}{9}

\bibitem{FED}
F.~Devvrit, A.~Krim-Yee, N.~Kumar, G.~MacGillivray, B.~Seamone, V.~Virgile, and
A.~Xu,
\emph{Fractional eternal domination: securely distributing resources across a
network},
Discuss. Math. Graph Theory \textbf{44} (2024), no.~4, 1395--1428.
\href{https://doi.org/10.7151/dmgt.2503}{doi:10.7151/dmgt.2503};
\href{https://arxiv.org/abs/2304.11795}{arXiv:2304.11795}.
Statements, the problem quoted in Section~1, and all verbatim definitions are
cited from the arXiv e-print, which has a single version.

\bibitem{KSV}
A.~Krim-Yee, B.~Seamone, and V.~Virgile,
\emph{Eternal domination on prisms of graphs},
Discrete Appl. Math. \textbf{283} (2020), 734--736.
\href{https://doi.org/10.1016/j.dam.2020.01.032}{doi:10.1016/j.dam.2020.01.032};
\href{https://arxiv.org/abs/1902.00799}{arXiv:1902.00799}.
Concerns the integer parameter on prisms and does not treat the fractional
game; listed because its title is close to the present subject.

\end{thebibliography}

\end{document}
